% r03_fill_K2a: proof of (K_0), part II: the Cameron-Martin transfer, the choice of the number of steps, effective actions, matching.

\subsubsection*{The Cameron--Martin transfer of the curvature source}

\begin{lemma}[Cameron--Martin transfer]\label{rot:lem:CM}
Let $\Lambda$ be a lattice of spacing $s=\eps_k$ with frame $(f_\alpha)$ descending to a torus $T$, $\gamma_\Lambda$ the normalised Gaussian photon measure of the action, with covariance $C_\Lambda=(-\Delta_\Lambda+\lmu^2)^{-1}\otimes\mathbf 1_3$ for the pairing $\ip uv_\Lambda=\sum_xs^3\sum_\alpha u_\alpha(x)v_\alpha(x)$, and write $\Zun_{\Lambda,T}(G,h)=\int d\gamma_\Lambda(A)\,e^{F_\Lambda(G)}W_\Lambda(A;h)$ with $W_\Lambda$ the matter integral, which depends on $A$ only through the bond phases. Then $F_\Lambda(G)=\ip A{\ell_\Lambda}_\Lambda$ with $\ell_{\Lambda,\beta}\defeq\sum_{\alpha\ne\beta}\partial^*_\alpha G_{\alpha\beta}$, $G_{\alpha\beta}(x)\defeq G_{c(x;\alpha\beta)}(f_\alpha,f_\beta)$, $\partial^*_\alpha u(x)\defeq s^{-1}(u(x-sf_\alpha)-u(x))$, and for real $z$
\[
\Zun_{\Lambda,T}(zG,h)=e^{z^2q_\Lambda(G)/2}\,\Zun_\Lambda[za_\Lambda](T;h),\qquad q_\Lambda(G)\defeq\ip{\ell_\Lambda}{C_\Lambda\ell_\Lambda}_\Lambda,\quad a_\Lambda\defeq C_\Lambda\ell_\Lambda,
\]
where $\Zun_\Lambda[b](T;h)\defeq\int d\gamma_\Lambda(A)\,W_\Lambda(A+b;h)$. Counterterms, Wick constant and wall are untouched. At $e=0$, $W_\Lambda$ does not depend on $A$, so $\langle e^{zF(G)}\rangle_{\fc,\Lambda}=e^{z^2q_\Lambda(G)/2}$ exactly.
\end{lemma}
\begin{proof}
$\sum_{\alpha<\beta}G_{\alpha\beta}(\partial_\alpha A_\beta-\partial_\beta A_\alpha)=\sum_{\alpha\ne\beta}G_{\alpha\beta}\partial_\alpha A_\beta$ by antisymmetry, and summation by parts on the torus gives $\ip A{\ell_\Lambda}$. Completing the square, $e^{z\ip A\ell}d\gamma(A)=e^{z^2\ip\ell{C\ell}/2}d\gamma(A-zC\ell)$; substitute $A\mapsto A+zC\ell$. Both sides are finite ($\lmu>0$, $\lambda>0$, finite lattice).
\end{proof}

\begin{lemma}[the Gaussian factor is frame-free]\label{rot:lem:q}
With $q(G)\defeq\int_T\ip{\dstar G}{(-\Delta+\lmu^2)^{-1}\dstar G}\,dx$: $|q_\Lambda(G)-q(G)|\le C(\lmu)s^2\|G\|_{H^3(T)}^2+C_N(G)s^N$ for every $N$ and every lattice $\Lambda=O(s\Z^3)$ descending to $T$.
\end{lemma}
\begin{proof}
In Fourier variables on $\Lambda/P$, with $k_\alpha\defeq k\cdot f_\alpha$, $\sigma_\alpha(k)\defeq\frac2s\sin\frac{sk_\alpha}2$ and $u_{\alpha\beta}(y)\defeq G_y(f_\alpha,f_\beta)$, Poisson summation on the coset of the plaquette centres gives $\hat\ell_{\Lambda,\beta}(k)=-\mathrm ie^{\mathrm isk_\beta/2}\sum_\alpha\sigma_\alpha(k)[\hat u_{\alpha\beta}(k)+r_{\alpha\beta}(k)]$ with alias terms $|r|\le C_Ns^N$ on the Brillouin zone, and $q_\Lambda=|T|^{-1}\sum_{k\in B_\Lambda\cap P^*}\sum_\beta|\hat\ell_{\Lambda,\beta}(k)|^2/(\sum_\gamma\sigma_\gamma(k)^2+\lmu^2)$. In the continuum, $q=|T|^{-1}\sum_{k\in P^*}(|k|^2+\lmu^2)^{-1}\sum_\beta|\sum_\alpha k_\alpha\hat u_{\alpha\beta}(k)|^2$, and $\sum_\beta|\sum_\alpha k_\alpha G(f_\alpha,f_\beta)|^2=\sum_\beta|G(k,f_\beta)|^2=|\iota_kG|^2$ is frame-free because $\sum_\alpha k_\alpha f_\alpha=k$. Finally $|\sigma_\alpha-k_\alpha|\le s^2|k|^3/24$, $|k|^2-s^2|k|^4/12\le\sum\sigma_\gamma^2\le|k|^2$, $\sum\sigma_\gamma^2\ge\frac4{\pi^2}|k|^2$ on $B_\Lambda$, and the continuum terms with $k\notin B_\Lambda$ sum to $O(s^N)$.
\end{proof}

\begin{lemma}[the shift field is regular]\label{rot:lem:shift}
Let $a^{\rm c}\defeq(-\Delta+\lmu^2)^{-1}\dstar G$, a smooth co-exact $1$-form on $T$. For every bond $(x,x+sf_\beta)$ of $\Lambda$, $|a_{\Lambda,\beta}(x)-(a^{\rm c}\cdot f_\beta)(x+\frac s2f_\beta)|\le C(G,\lmu)s^2$, and the same for lattice differences of any fixed order. In King's unit scaling at step $r$ the background $za_\Lambda$ is regular in the sense of \cite[Def.~3.2]{K1} for $\delta_r\le\delta_0(z,G,e)$, and so is $A^{(r)}+za_\Lambda$ whenever $A^{(r)}$ is.
\end{lemma}
\begin{proof}
$\hat a_{\Lambda,\beta}=(\sum_\gamma\sigma_\gamma^2+\lmu^2)^{-1}\hat\ell_{\Lambda,\beta}$; compare coefficient by coefficient with the sampled continuum field as in \Cref{rot:lem:q}. In unit scaling ($A_u(x_u)=\delta_r^{1/2}A(\delta_rx_u)$, $e_u=e\delta_r^{1/2}$, \cite[(2.20)]{K1}) a derivative of $za_\Lambda$ is $\le|z|\delta_r^{3/2}(\|\nabla a^{\rm c}\|_\infty+Cs^2)$, while \cite[Def.~3.2]{K1} allows $Ce^{\beta-1}\delta_r^{(\beta-1)/2}$ with a small $\beta>0$.
\end{proof}

\begin{remark}[what the transfer does for the comparison]\label{rot:rem:CMrg}
After \Cref{rot:lem:CM} the curvature source has disappeared from the photon: the model is King's with a scalar source and a fixed smooth external background $za_\Lambda$ in the matter coupling only. Running the matter averages in the total field, the background enters exactly where King's coarse background $A^{(r)}$ enters, through $A^{(r)}+za_\Lambda$, regular and small at every scale below $\delta_0(z,G)$ (\Cref{rot:lem:shift}); this is the content of the identities \cite[(3.82)--(3.87)]{K1}, which express the effective action at a background as the action at zero background with the background inserted in the lines and vertices. The two backgrounds $a_{\Xk}$ (axis bonds) and $a_{\Yk}$ (rotated bonds) are samples of one continuum $1$-form $a^{\rm c}$, so replacing one by the other at matched points is the same step as King's replacement of $g_{\mu\nu}(x')$ by $g_{\mu\nu}(x)$ \cite[p.~665]{K1}. The comparison of the effective actions below is therefore carried out at ``$h=0$ with background'', and no frame-paired source vertex \cite[(3.20)]{K1} appears.
\end{remark}

\begin{lemma}[lattice sums]\label{rot:lem:lsum}
For $h\in C^\infty(T)$, every lattice $\Lambda=O(s\Z^3)\supset P$ and every $N$: $|\sS_\Lambda(h)-\int_Th|\le C_N(h)s^N$, with $C_N(h)$ independent of $O$.
\end{lemma}
\begin{proof}
$\sS_\Lambda(h)=\sum_{k\in\Lambda^*}\hat h(k)$ by Poisson summation on $\Lambda/P$ ($P\subset\Lambda$ gives $\Lambda^*\subset P^*$); the term $k=0$ is $\int h$; every $k\in\Lambda^*\setminus0$ has $|k|\ge2\pi/s$; $|\hat h(k)|\le\|\Delta^Mh\|_{L^1}|k|^{-2M}$ and $\sum_{k\in\Lambda^*\setminus0}|k|^{-2M}\le C_Ms^{2M}$ for $2M>3$.
\end{proof}

\subsubsection*{The number of steps, the effective actions, the matching}

\begin{lemma}[the choice of $r$]\label{rot:lem:r}
Let $\gamma_*,\beta_*,\sigma>0$ be the exponents of \cite[Thm.~3.4]{K1}, $\beta_1$ the field-norm exponent of step (5) of the proof of \Cref{rot:thm:A}, $\theta_1$ the gain exponent of \Cref{rot:thm:H1S}, $\gamma\defeq\min(\gamma_*,\theta_1,\frac12)$ and $\beta\defeq\max(\beta_*,\beta_1+\frac12,5)$. Fix $0<\vr<\min\{\gamma/(\gamma+\beta),\frac19\}$ and put $k-r(k)\defeq\lfloor\vr k\rfloor$. Then, as $k\to\infty$: $\delta_r\to0$; $5^{-\gamma r}\delta_r^{-\beta}(1+k)^p\to0$ for every $p$; $5^{3k-5r}\to0$; $\delta_r^{-2}(\eps_k/\delta_r)^{1/4}\log(1/\eps_k)^2\to0$; and $\eps_k^{-q}e^{-c\delta_r^{-1}}\to0$ for all $q,c>0$. Moreover $p(\delta_r)^2\ge b_0^2(1+\vr k\ln5-\ln5)^{2\mathrm p}$ grows at least quadratically in $k$, where $p(\delta)=b_0(1+\log\delta^{-1})^{\mathrm p}$, $\mathrm p\ge1$, is King's small-field function \cite[(3.1)]{K1}.
\end{lemma}
\begin{proof}
$\delta_r=5^{-\lfloor\vr k\rfloor}$. The second quantity is $5^{-\gamma k+(\gamma+\beta)\lfloor\vr k\rfloor}(1+k)^p$; the third $\le5^{-2k+5\vr k+5}$; the fourth $\le5^{2\vr k+2-(1-\vr)k/4}(k\ln5)^2$, and $2\vr<(1-\vr)/4$ iff $\vr<\frac19$; the fifth because $\delta_r^{-1}\ge5^{\vr k-1}$ grows exponentially. The last claim is \cite[(3.1)]{K1} with $\mathrm p\ge1$. King's own choice \cite[(3.12)]{K1}, $\vr=\gamma/(2\beta+\gamma)$, satisfies every condition.
\end{proof}

\begin{definition}[effective actions; vertices]\label{rot:def:eff}
For the standard flow the $r$-step effective action is \cite[(3.14)]{K1} rewritten as \cite[(3.90)]{K1}:
\[
S^{(r)}=\tfrac12\ip{A_r}{\Delta^{(r)}A_r}+\ln N_r+\tfrac12\ip{\phi_r}{\Delta^{(r)}(A^{(r)})\phi_r}+\ln N_r(0)-\ln\big[Z_r(A^{(r)})Z_r(0)^{-1}\big]+P^{(r)}(A_r,\phi_r;s),
\]
where $P^{(r)}$ is the Taylor polynomial of order $\bar n=13$ in the running couplings of King's $E'_r$ \cite[(3.17)]{K1}, a finite sum of connected graphs with the vertices \cite[(3.18)--(3.26)]{K1}. The hybrid $\hs^{(r)}$ is the same expression for the hybrid flow. A curvature source enters only through the Cameron--Martin background (\Cref{rot:rem:CMrg}); a wall enters through the source vertex \cite[(3.21)]{K1}, the \emph{$h$-vertex}, $\theta(5^r\eps)^2\sum_x\eta^3h(x)\sum_i[(\phi'_i+\phi^{(r)}_i)(x)^2-(-\Delta^{\eta,\Omega}+m^2(5^r\eps)^2)^{-1}(x,x)]$, whose subtraction is in \cite{K1} the coincident torus propagator at the running scale and here, in the free-boundary convention which both flows inherit from the fine action by exact integration, the corresponding infinite-lattice coincident propagator, the same number for both lattices because $\Xk$ and $\Yk$ are isometric; the two readings differ by the image sum $\sum_{\gamma\in P\setminus0}$ of the running propagator, a $c$-number times $\sS(h)$, which is part of the bridge \Cref{rot:prop:CT}. Everything below is written in the free-boundary reading.
\end{definition}

\begin{definition}[matching; slices]\label{rot:def:match}
In the unit scale of level $r$ let $\Gamma_r=R^{-1}\Z^3$ (the coarse lattice), $\Gamma_1=R^{-1}(\eta\Z^3)$ ($=Y_{k-1}$ rescaled), $\Lambda_X=\eta'\Z^3$ and $\Lambda_Y=R^{-1}(\eta'\Z^3)$. For $x\in\Gamma_1$ the first-step cells are the tile $x+\eta'B_0\subset\Lambda_X$ and the centred cube $Q_Y(x)\subset\Lambda_Y$, $125$ sites each. Fix bijections $\tau_x:x+\eta'B_0\to Q_Y(x)$; they glue to a bijection $\tau:\Lambda_X\to\Lambda_Y$, and on every admissible torus to $\Xk/P\to\Yk/P$, with $|\tau x'-x'|\le c_\tau\eta'$, $c_\tau=\sqrt{13}+2\sqrt3<7.1$. Points $x''=\tau x'$ are \emph{matched}. For a flow $F\in\{X,Y\}$ let $S^F_0\defeq G^F_1$ be the first fluctuation covariance and $S^F_j\defeq H^F_jC^F_jH^{F*}_j$ the $(j+1)$-st, lifted to the fine lattice; then $G^F_r=\sum_{j<r}S^F_j$ exactly, and on the infinite lattice $C^F=G^F_n+\Gamma^F_n$ with $\Gamma^F_n\defeq H^F_n(A^F_n)^{-1}H^{F*}_n$ (Woodbury). Slice $j$ has length scale $\ell_j\defeq5^j\eta'$. The \emph{coarse slices} $j\ge r$ are the further steps of the same flow; for $F=X$ these are standard rotated steps on $Y_{k-r},Y_{k-r-1},\dots$
\end{definition}

\begin{lemma}[vertex matching; frame-free contractions]\label{rot:lem:vert}
\textup{(i)} $\sum_{y\in\Lambda_Y^n}\eta'^{3n}g(y)=\sum_{x\in\Lambda_X^n}\eta'^{3n}g(\tau x)$: every $Y$-graph is, exactly, a sum over $X$-vertices with $\tau$-transported kernels. \textup{(ii)} A smooth vertex function $w$ (couplings, the Cameron--Martin background, partitions of unity) sampled at $x'$ and at $\tau x'$ differs by $\le c_\tau\eta'\|\nabla w\|_\infty$. \textup{(iii)} Vector indices are contracted by the Euclidean inner product and derivative lines are compared as vectors; the current vertex $eA\cdot j$ (\cite[(3.22)]{K1}, $m=1$) and the seagull $e^2|A|^2|\phi|^2$ (the $m=2$ member of \cite[(3.23)]{K1}) are frame-free; \cite[(3.22)]{K1} with $m=2$ is the bond-diagonal derivative vertex $\frac{e^2\eta}2\sum_bA_b^2D_b\phi\cdot U\phi(b_+)$ and is not frame-free; the vertices with $m\ge3$ in (3.22) and $m\ge4$ in (3.23) carry explicit factors $\eta'^{m-2}$ or more and are bounded per flow.
\end{lemma}
\begin{proof}
(i) $\tau$ is a bijection of $\Lambda_X$ onto $\Lambda_Y$ and both lattices have the cell volume $\eta'^3$ per site, so $y=\tau x$ is an exact change of the summation variable in each factor; on an admissible torus the same holds for the bijection $\Xk/P\to\Yk/P$. (ii) The mean value theorem along the segment from $x'$ to $\tau x'$. (iii) In the unit scale of level $r$ the hopping term on the fine lattice of spacing $\eta'$ is $\sum_b\eta'\big[|\phi(b_-)|^2+|\phi(b_+)|^2-2\Re\big(\bar\phi(b_+)e^{\mathrm ie_r\eta'A_b}\phi(b_-)\big)\big]$, the bond integral of the unit-scaled field being $\eta'A_b$ \cite[(2.20)]{K1} and $e_r=e(5^r\eps)^{1/2}$ the running charge; its expansion in powers of $A_b$ carries $(e_r\eta')^m/m!$ at order $m$, one explicit $\eta'$ per power of $A$, which is how \cite[(3.22)--(3.23)]{K1} are organised. At $m=1$ the vertex is $\sum_x\eta'^3e_rA(x)\cdot j(x)$ with $j=\sum_\alpha j_\alpha f_\alpha$ the lattice current, $j_\alpha(x)=2\,\Im(\bar\phi(x+\eta'f_\alpha)\phi(x))/\eta'$; at $m=2$ the member of (3.23) is $\sum_x\eta'^3e_r^2|A(x)|^2|\phi(x)|^2$. The pairing $A\cdot j$ and the norm $|A|^2$ are the same functions of the vectors in every orthonormal frame; comparing the $X$-version with the $\tau$-transported $Y$-version therefore compares the vectors $A$, $j$, $\phi$ at matched points only. The member of (3.22) with $m=2$ contracts $A\otimes A\otimes D\phi$ with the rank-three tensor $\sum_\alpha f_\alpha\otimes f_\alpha\otimes f_\alpha$, which is not $O(3)$-invariant (there is no non-zero isotropic symmetric tensor of rank three), so the two versions differ at order one on the same vectors; its tadpole is handled by the summation by parts of \Cref{rot:lem:CK4}. The vertices with $m\ge3$, $m\ge4$ carry $\eta'^{m-2}$ relative to the two order-$e_r^2$ vertices, and each flow's own line bounds \cite[(3.63)--(3.65)]{K1} bound them without any comparison.
\end{proof}

\begin{lemma}[counterterms take the same value along both flows]\label{rot:lem:ctsame}
Let $g$ be a counterterm graph of $\dm$ or $E_1$ (free boundary conditions), an infinite-lattice graph with full lines $C$ localised at one vertex. Expanding each of its lines in the slices of flow $F$ \cite[(3.76)]{K1}, the sum over all slice assignments equals $g$ for both flows. Mixed assignments take other values; they occur only inside a telescoping sum whose endpoints are the pure assignments.
\end{lemma}
\begin{proof}
$\sum_jS^F_j=C^F$ pointwise; $C^Y$ is the image of $C^X$ under the isometry $R^{-1}$, and the localised vertex sum of an infinite-lattice graph is invariant under that isometry.
\end{proof}
